\section{Qserv catalogue database}\label{sec:upqserv}

This replaces the Qserv rows of \tabref{tab:storageSizingOps},
\tabref{tab:datasetSizingOps}, \tabref{tab:storageFloorOps} and
\tabref{tab:computeSizingOps}.

\subsection{Deployed hardware}

Worker nodes, as confirmed by the Qserv team in July 2026:

\begin{itemize}
\item CPU: one AMD EPYC 7543P (32 cores, 64 hardware threads, SMT enabled).
\item RAM: 256\,GB.
\item Storage: node-local NVMe, shared-nothing, with no network storage in
      the query path. 12 of 24 slots are populated, giving 32\,TB raw per
      node. ZFS compression gives about a factor of two, so about 64\,TB is
      usable per node. Populating the remaining slots would double raw
      capacity.
\item Internal network: 100\,GbE.
\item Replication: three-fold on production.
\end{itemize}

As of July 2026 the named deployments were:
\begin{itemize}
\item a 35-node Kubernetes production cluster, serving DP2;
\item a 35-node Kubernetes integration cluster;
\item a 6-node Docker integration cluster.
\end{itemize}
The model counts \smQsNodesHave{} Qserv-class nodes in service, from the March
2026 inventory. That is consistent with the Qserv team's September 2026
description of about 90 nodes in two clusters. Most USDF Kubernetes nodes share
this hardware. Which cluster will serve DR1 is to be confirmed.

\subsection{The measured DP2 catalogue}

The data-preview column is measured, not modelled. On the USDF deployment the
DP2 catalogue occupies \smQsDPtwoLogicalTB{}\,TB for one copy of each chunk:
27.0\,TB of chunk tables plus 6.4\,TB of overlap tables, of which
\smQsDPtwoIndexTB{}\,TB are indexes. With every replica it is 66.9\,TB on disk,
exactly twice the single-copy figure. The deployment measured therefore runs
two-fold replication. Production is to run three-fold, so the model
provisions \smQsDPtwoDiskTB{}\,TB for the preview. That figure is a
provisioning requirement, not a measurement.

\subsection{How data-release catalogues are sized}

The single-copy logical volume of each data release is computed from
\citeds{LDM-141} row counts and row sizes, as in \tabref{tab:datasetSizingOps}:
\[
  V_{\mathrm{logical}} =
  \bigl(N_{\mathrm{obj}}\,(r_{\mathrm{obj}} + r_{\mathrm{obj,extra}})
        + N_{\mathrm{src}}\,r_{\mathrm{src}}\bigr)\,g^{k}
  + N_{\mathrm{fs}}\,r_{\mathrm{fs}} ,
\]
where $g$ is one plus the row-size growth per release and $k$ counts releases
since DR1. The volume on disk is
\[
  V_{\mathrm{disk}} = V_{\mathrm{logical}} \times (1 + o) \times R ,
\]
with overlap uplift $o = \smQsOverlap$ and replication $R = \smQsRep$.

\smpara{Compression is applied once, on the capacity side.}
ZFS compression is already in the usable capacity per node:
$32\,\mathrm{TB} \times 2 \times 0.675$ usable $= 43.2$\,TB. The catalogue
volume is therefore carried uncompressed.

\smpara{Row sizes grow between releases.}
As in \tabref{tab:storageSizingOps}, row sizes widen as measured columns are
added. The construction-era values imply 3.02\%, 3.00\% and 3.12\% per step for
Object, Object\_Extra and Source. The model uses \smRowGrowth{} per release,
reaching a factor of \smQsDRnineInfl{} by DR9. Forced sources are held flat,
because that schema is fixed.

\smpara{Forced sources include those on DIAObjects.}
The LDM-141 forced-source count covers forced photometry on both Objects and
DIAObjects, which DP2 holds in two tables (ForcedSource and
ForcedSourceOnDIAObject), so the second table is not added separately. This
assumes both have roughly the 41-byte row size.

\smpara{Overlap.}
The uplift of \smQsOverlap{} is measured on DP2: overlap bytes over chunk bytes,
for the whole catalogue. The partitioning uses 340 stripes of 3 sub-stripes,
giving sub-chunks $180^\circ / 340 / 3 = 10.59$ arcmin high. That is identical
to the UK Data Facility's configuration. With a 1-arcmin overlap, the
geometric overlap is about 41\% by area.

The measured aggregate is lower because only director tables carry overlap;
child tables such as ForcedSource carry none. As forced sources accumulate,
the director share of the catalogue falls from about half at DR1 to about a
sixth at DR9. A flat aggregate uplift therefore overstates later releases, by
up to about 13\% at DR9.

\smpara{Chunk count follows sky area, not data volume.}
Partitioning is fixed across a database family, so chunk count grows with sky
coverage. DP2's 57{,}680 chunks already cover an estimated 11{,}700\,deg$^2$,
so chunk growth from DP2 to DR1 is modest.

\subsection{Calibration}

Two checks sit beside the projection in \tabref{tab:smQserv}:

\begin{itemize}
\item An independent planning estimate puts DR1 at
      \smQsPlanLowPB--\smQsPlanHighPB\,PB for one copy. The model gives
      \smQsDRonePerCopyPB\,PB, or \smQsDRoneVsPlan{} of that range's midpoint.
      The estimate's source has not been traced, and the gap is recorded
      rather than adjusted away.
\item The modelled DR1 catalogue is more than forty times the measured DP2
      catalogue, for about seven times DP2's visits. Some of that is
      expected, because DR1 covers more sky and forced photometry grows with
      depth. The rest suggests a mismatch between LDM-141 row sizes and the
      schema actually deployed. Per-table row counts for DP2 are needed to
      settle it. Until then, read the data-release columns as an upper bound on
      the LDM-141 basis.
\end{itemize}

\subsection{Node projection}

Usable data per node is \emph{assumed}:
\begin{itemize}
\item 43.2\,TB in years 1 and 2;
\item 86.4\,TB from year 3, with all 24 NVMe slots populated;
\item 172.8\,TB from year 7, with denser drives.
\end{itemize}
Qserv holds a sliding window of three releases: two served, one in
preparation, as in the construction-era model.

On these assumptions year 2, when the first data release is loaded, needs
\smQsDRoneNodes{} nodes against \smQsNodesHave{} in service, a gap of
\smQsDRoneGap{}. That gap depends entirely on when the empty slots are filled.
If they are populated before the first release is loaded, the same
\smQsDRoneWindowTB{}\,TB needs about half as many nodes. By year 10 the
projection reaches \smQsDRnineNodes{} nodes. The fall in node count in year 7
comes from the drive-density step, not from falling demand.
\textbf{Populating the empty NVMe slots before the first data release is loaded
is the most effective Qserv action available.}

% GENERATED by generate_model.py from lsst/rubin-sizing-model @ v2026.10.0. Do not edit by hand.
% Regenerate: git clone --branch v2026.10.0 https://github.com/lsst/rubin-sizing-model ; then, inside the clone, uv run generate_model.py --emit-tex <this directory>

\small \begin{longtable}{|p{0.55\textwidth}|r|l|}
\caption{Qserv inputs: deployment, replication, overlap, LDM-141 row sizes and their growth, and the measured data-preview catalogue. \label{tab:smQservInputs}}\\
\hline
\textbf{Metric} & \textbf{Value} & \textbf{Unit} \\
\hline\endfirsthead
\hline
\textbf{Metric} & \textbf{Value} & \textbf{Unit} \\
\hline\endhead
Qserv nodes deployed & 92 & nodes \\ \hline
Qserv replication factor & 3 & x \\ \hline
Overlap uplift on chunk bytes & 24\% & fraction \\ \hline
Qserv retention window & 3 & releases \\ \hline
Object row size & 1,896 & bytes \\ \hline
Object extra row size & 21,005 & bytes \\ \hline
Source row size & 467 & bytes \\ \hline
Forced-source row size & 41 & bytes \\ \hline
Row-size growth per release & 3\% & fraction \\ \hline
LDM row size to stored-bytes factor & 1.000 & ratio \\ \hline
Measured data-preview catalogue (logical) & 33.45 & TB \\ \hline
Measured data-preview indexes (logical) & 3.24 & TB \\ \hline
Independent DR1 catalogue estimate, low & 2.0 & PB \\ \hline
Independent DR1 catalogue estimate, high & 3.0 & PB \\ \hline
\end{longtable} \normalsize


\begin{landscape}
% GENERATED by generate_model.py from lsst/rubin-sizing-model @ v2026.10.0. Do not edit by hand.
% Regenerate: git clone --branch v2026.10.0 https://github.com/lsst/rubin-sizing-model ; then, inside the clone, uv run generate_model.py --emit-tex <this directory>

\tiny \begin{longtable}{|p{0.22\textwidth}|l|r|r|r|r|r|r|r|r|r|r|}
\caption{Qserv catalogue size and node projection. The first column is the measured data-preview catalogue; later columns are modelled from LDM-141 row sizes. \label{tab:smQserv}}\\
\hline
\textbf{Metric} & \textbf{Unit} & \textbf{Year 1 (DP2)} & \textbf{Year 2} & \textbf{Year 3} & \textbf{Year 4} & \textbf{Year 5} & \textbf{Year 6} & \textbf{Year 7} & \textbf{Year 8} & \textbf{Year 9} & \textbf{Year 10} \\
 &  & FY2026 & FY2027 & FY2028 & FY2029 & FY2030 & FY2031 & FY2032 & FY2033 & FY2034 & FY2035 \\
\hline\endfirsthead
\hline
\textbf{Metric} & \textbf{Unit} & \textbf{Year 1 (DP2)} & \textbf{Year 2} & \textbf{Year 3} & \textbf{Year 4} & \textbf{Year 5} & \textbf{Year 6} & \textbf{Year 7} & \textbf{Year 8} & \textbf{Year 9} & \textbf{Year 10} \\
 &  & FY2026 & FY2027 & FY2028 & FY2029 & FY2030 & FY2031 & FY2032 & FY2033 & FY2034 & FY2035 \\
\hline\endhead
Objects & billions &  & 27.5 & 32.5 & 35.7 & 38.2 & 40.3 & 42.2 & 43.8 & 45.3 & 46.4 \\ \hline
Sources & billions &  & 901 & 1,800 & 2,700 & 3,600 & 4,510 & 5,410 & 6,310 & 7,210 & 8,110 \\ \hline
Forced sources & billions &  & 2,910 & 6,870 & 11,300 & 16,100 & 21,300 & 26,700 & 32,400 & 38,300 & 44,100 \\ \hline
\quad row-size inflation vs DR1 & x &  & 1.000 & 1.030 & 1.061 & 1.093 & 1.126 & 1.159 & 1.194 & 1.230 & 1.267 \\ \hline
Catalogue, logical uncompressed, one copy & TB &  & 1,170 & 1,914 & 2,668 & 3,453 & 4,283 & 5,144 & 6,045 & 6,987 & 7,952 \\ \hline
Object table, replicated & TB &  & 156 & 185 & 203 & 217 & 229 & 240 & 249 & 258 & 264 \\ \hline
Fast (low-latency) tier: APDB + Object table & TB & 17 & 190 & 235 & 269 & 300 & 329 & 356 & 382 & 407 & 430 \\ \hline
Worker database on disk, this release (replicated, with overlap) & TB & 100 & 4,341 & 7,103 & 9,902 & 12,815 & 15,893 & 19,089 & 22,432 & 25,930 & 29,510 \\ \hline
Worker database, retention window & TB & 100 & 4,442 & 11,545 & 21,347 & 29,820 & 38,609 & 47,796 & 57,414 & 67,451 & 77,871 \\ \hline
Worker database per copy & TB & 33 & 1,447 & 2,368 & 3,301 & 4,272 & 5,298 & 6,363 & 7,477 & 8,643 & 9,837 \\ \hline
As a multiple of the measured data preview & x & 1.0 & 43.3 & 70.8 & 98.7 & 127.7 & 158.4 & 190.2 & 223.5 & 258.4 & 294.1 \\ \hline
Visits as a multiple of the data preview & x & 1.0 & 7.0 & 14.0 & 21.0 & 28.0 & 35.0 & 42.0 & 49.0 & 56.0 & 63.0 \\ \hline
Independent DR1 estimate, per copy (midpoint) & TB &  & 2,500 &  &  &  &  &  &  &  &  \\ \hline
Model as a multiple of that estimate & x &  & 0.58 &  &  &  &  &  &  &  &  \\ \hline
Data per node (drive-density ramp) & TB/node & 43.2 & 43.2 & 86.4 & 86.4 & 86.4 & 86.4 & 172.8 & 172.8 & 172.8 & 172.8 \\ \hline
QSERV NODES NEEDED & nodes & 3 & 103 & 134 & 248 & 346 & 447 & 277 & 333 & 391 & 451 \\ \hline
Qserv nodes deployed today & nodes & 92 & 92 & 92 & 92 & 92 & 92 & 92 & 92 & 92 & 92 \\ \hline
Surplus / (gap) & nodes & 89 & -11 & -42 & -156 & -254 & -355 & -185 & -241 & -299 & -359 \\ \hline
\end{longtable} \normalsize

\end{landscape}
